% !PRINTED 609
\weberpage{609}
Every integer $\eta$ of the field $\mathfrak{K}$ can (Bd.\,I, § 151) be represented in
the form:
\[
\eta = \frac{\varphi(\xi)}{f'(\xi)},
\]
where $\varphi(\xi)$ has integral coefficients in $\Omega$. One can put this in the
form:
\weq{9}{\gamma\eta = \gamma_0 + \gamma_1\xi + \gamma_2\xi^2 + \cdots + \gamma_{n-1}\xi^{n-1},}
where $\gamma, \gamma_0, \gamma_1, \ldots, \gamma_{n-1}$ are integers in $\Omega$, and such that the
first, $\gamma$, is a divisor of the discriminant of $f$. The prime factors of
$\gamma$ we exclude by means of the excludent.

Let

$\mathfrak{P}$ be a prime ideal of the first degree in $\mathfrak{K}$,

$\mathfrak{p}_1$ the prime ideal of the first degree in $\Omega$ that is divisible by $\mathfrak{P}$,

$p$ the norm of $\mathfrak{p}_1$ and of $\mathfrak{P}$;

then for every integer $\alpha$ in $\Omega$:
\[
\alpha^p \equiv \alpha \ (\mathrm{mod}\ \mathfrak{p}_1)
\]
and
\weq{10}{\xi^p \equiv \xi \ (\mathrm{mod}\ \mathfrak{P}),}
hence, by (9), for every integer $\eta$ in $\mathfrak{K}$:
\[
\eta^p \equiv \eta \ (\mathrm{mod}\ \mathfrak{P}).
\]

9. \gesperrt{The necessary and sufficient condition for $\mathfrak{K}(\xi)$ to be a class field therefore consists in this, that for every prime ideal $\mathfrak{P}$ dividing a $\mathfrak{p}_1$, and only for this one, the congruence (10) is satisfied.}

Let now
\[
\mathfrak{K} = \Omega(\xi), \quad \mathfrak{K}' = \Omega(\xi')
\]
be two class fields of the same group $A$. First of all, it readily follows that they must be \gesperrt{of the same degree}. For the condition (5) must be satisfied for both. If, then, $n$ and $n'$ are the degrees, $(s-1)^{n'-n}$ must be finite and different from zero for $s=1$, which is only possible if $n=n'$.

By composition of $\mathfrak{K}$ and $\mathfrak{K}'$ we derive a third field:
\[
\mathfrak{K}'' = \Omega(\xi, \xi'),
\]
of which we shall show that it too is a class field.
